\tzpanchor{0349}
\tzpentryheader{0349}{Tripartite Information Retrieval Algorithm}

\begingroup\small
\begingroup\scriptsize
\setlength{\tabcolsep}{4pt}
\renewcommand{\arraystretch}{0.92}
\noindent\begin{tabularx}{\linewidth}{@{}p{0.24\linewidth}p{0.36\linewidth}X@{}}
\textbf{UUID} & \multicolumn{2}{l@{}}{\tzpcode{fda91122-0052-5919-b534-1093ebfc7e46}}\\
\textbf{PiID} & \tzpcode{6789259036001} & \textbf{TZPi}\quad \tzpcode{9}\quad \textbf{PiPE}\quad \tzpcode{356}\\
\textbf{RTEID} & \textbf{Poloidal $\theta$}\quad \tzpcode{1355.2445 rad} & \textbf{Toroidal $\phi$}\quad \tzpcode{02.2330 rad}\\
\textbf{TZPID} & \tzpcode{ID0349} & \textbf{Node Dependencies}\quad \tzpidref{0348}\\
\textbf{Registry} & \tzpcode{registry\_id\_0349} & \textbf{Scale}\quad informational\\
\textbf{LOC Classification} & QA76.9 information and coding & QA614 topological data analysis\\
\textbf{arXiv Classification} & Primary: cs.IT & Cross-list: math.AT, math-ph\\
\textbf{Primary Support} & Carrier Map & Composable Action\\
\textbf{Closure Support} & Normalized Output & \\
\end{tabularx}\par
\endgroup
\endgroup

\tzpsection{Canonical Statement}
To retrieve information across the tripartite channels: (1) measure sequences of DAANSphere addresses corresponding to mediator channel events; (2) decode these sequences via the holographic braille dictionary to identify topological features on the null surface; (3) reconstruct classical information by projecting the identified features onto the temporal channel, subject to decoherence and holographic bounds.

\tzpsection{Canonical Equation}
\[
Q=(q_{A},q_{B},q_{C})
\]
\[
\mathcal{R}(Q)=\operatorname*{arg\,max}_{x}\sum_{j\in\{A,B,C\}} I(x:q_{j})
\]

\begin{minipage}[t]{0.49\linewidth}
\tzpsection{Canonical Symbol Key}
\begingroup
\small
\raggedright
\textbf{Source equation symbols.}\par
\begin{itemize}[leftmargin=1.35em,itemsep=1pt,topsep=1pt,parsep=0pt]
    \item $Q$: source equation symbol.
    \item $q_{A}$: source equation symbol.
    \item $A$: source equation symbol.
    \item $q_{B}$: source equation symbol.
    \item $B$: magnetic field magnitude.
    \item $q_{C}$: source equation symbol.
    \item $C$: speed of light or relativistic conversion factor.
    \item $\mathcal{R}(Q)$: global radius, boundary radius, or topological scale.
    \item $I$: source equation symbol.
    \item $q_{j}$: source equation symbol.
\end{itemize}
\endgroup
\end{minipage}\hfill
\begin{minipage}[t]{0.47\linewidth}
\tzpsection{Wolfram Form}
\begingroup
\scriptsize
\raggedright
\textbf{Source Wolfram model.}\par
\begin{Verbatim}[fontsize=\tiny,breaklines=true,breakanywhere=true]
(* TZPID ID0349 generated source Wolfram model *)
ClearAll[K0349, Cr0349, T, R, A, W, I, N];
eq0349$1 = HoldForm[Q == Tuple[qA, qB, qC]];
eq0349$2 = HoldForm[R[Q] == ArgMax[Sum[I[x, q[j]], {j, {A, B, C}}], x]];
eq0349$3 = HoldForm[x^ast == R[Q]];

K0349 = HoldForm[{eq0349$1, eq0349$2, eq0349$3}];

N = "normalized TRAWIN closure object for Tripartite Information Retrieval Algorithm";
I = "Normalized Output integration/comparison layer";
W = "Composable Action wave or operator carrier";
A = "Carrier Map admissible action/amplitude condition";
R = "Composable Action rotation/curl preservation layer";
T = "Carrier Map local source condition";

Cr0349[expr_] := HoldForm[N[I[W[A[R[T[expr]]]]]]];

Result0349 = Cr0349[K0349];
\end{Verbatim}
\endgroup
\end{minipage}

\par\vspace{0.45\baselineskip}
\noindent\begin{minipage}[t]{\linewidth}
\begingroup
\small
\raggedright
\textbf{TRAWIN Operator Key.}\par
\noindent\textbf{Time/localization operator:} $\mathsf{T}\!\left[\mathrm{local}\right]$ Carrier Map local source condition.\par
\noindent\textbf{Rotation/curl operator:} $\mathsf{R}\!\left[\nabla\times\right]$ Composable Action rotation/curl preservation layer.\par
\noindent\textbf{Action/amplitude operator:} $\mathsf{A}\!\left[\mathrm{admissible}\right]$ Carrier Map admissible action/amplitude condition.\par
\noindent\textbf{Wave/d'Alembertian operator:} $\mathsf{W}\!\left[\Box\right]$ Composable Action wave or operator carrier.\par
\noindent\textbf{Integration operator:} $\mathsf{I}\!\left[\int\right]$ Normalized Output integration/comparison layer.\par
\noindent\textbf{Normalization operator:} $\mathsf{N}\!\left[\mathrm{normalized}\right]$ normalized TRAWIN closure object for Tripartite Information Retrieval Algorithm.\par
\endgroup
\end{minipage}

\clearpage
\noindent\textbf{[ID 0349] Tripartite Information Retrieval Algorithm}\par
\vspace{0.35em}
\tzpsection{Derivation Layer}
\paragraph{Primary Equation.}
\begin{equation}
\left\langle\phi^2\right\rangle
=
\int_{0}^{\omega_{\mathrm{cut}}}
\frac{\hbar}{2\omega}
\eta(\omega)
\,\mathrm{d}\omega .
\end{equation}

\paragraph{Primary Derivation.}
For a quantum field mode of frequency $\omega$, the vacuum fluctuation contribution scales as:

\begin{equation}
\frac{\hbar}{2\omega}.
\end{equation}

The spectral density index $\eta(\omega)$ weights the contribution of each frequency mode:

\begin{equation}
\frac{\hbar}{2\omega}\eta(\omega).
\end{equation}

Integrating over allowed frequencies from $0$ to the cutoff $\omega_{\mathrm{cut}}$ gives the total zero-point fluctuation variance:

\begin{equation}
\boxed{
\left\langle\phi^2\right\rangle
=
\int_{0}^{\omega_{\mathrm{cut}}}
\frac{\hbar}{2\omega}
\eta(\omega)
\,\mathrm{d}\omega .
}
\end{equation}

\paragraph{Interpretation.}
Zero-point fluctuations are represented as the weighted variance of field amplitudes over the allowed frequency spectrum.

\par\vspace{0.35em}
\noindent\textbf{Derivable Consequences.}\quad
\begingroup
\footnotesize
\raggedright
The canonical equation anchors Tripartite Information Retrieval Algorithm by connecting the source condition to the derivation layer and proof certificate.
\par\endgroup

\tzpsection{Equation Support}
\noindent\textbf{1. Carrier Map}\par
\[
Q=(q_{A},q_{B},q_{C})
\]
\noindent This names the entry as an actual map between carrier domains rather than a loose metaphor.\par
\noindent\textbf{2. Composable Action}\par
\[
\mathcal{R}(Q)=\operatorname*{arg\,max}_{x}\sum_{j\in\{A,B,C\}} I(x:q_{j})
\]
\noindent This inserts the map into the ordered TRAWIN chain and makes composition explicit.\par
\noindent\textbf{3. Normalized Output}\par
\[
x^{\ast}=\mathcal{R}(Q)
\]
\noindent This closes the entry by requiring that the output land in the admissible target class.\par

\clearpage
\noindent\textbf{[ID 0349] Tripartite Information Retrieval Algorithm}\par
\vspace{0.35em}
\tzpsection{Dictionary}
\begingroup\small
\tzpsection{Definition}
Tripartite Information Retrieval Algorithm is a registry definition in Field Theory and Action Principles with secondary pressure from Manifold and Metric Dynamics.

% BEGIN TZP CONTEXTUAL MEANING
\tzpsection{Contextual Meaning}
\begingroup\footnotesize\setlength{\emergencystretch}{3em}\sloppy
\textbf{Formal statement:} To retrieve information across the tripartite channels: (1) measure sequences of DAANSphere addresses corresponding to mediator channel events; (2) decode these sequences via the holographic braille dictionary to identify topological features on the null surface; (3) reconstruct classical information by projecting the identified features onto the temporal channel, subject to decoherence and holographic bounds. \textbf{Contextual meaning:} The entry reads Tripartite Information Retrieval Algorithm as a controlled technical headword whose equation layer, proof route, and registry role are interpreted together. \textbf{Derivable consequences:} The entry converts a registry title into a transport map that can be reused by the surrounding transfer arc. \textbf{Lemma support:} Carrier Map; Composable Action; Normalized Output.\par
\endgroup
% END TZP CONTEXTUAL MEANING

\tzpsection{Technical Interpretation}
Zero-point fluctuations are represented as the weighted variance of field amplitudes over the allowed frequency spectrum.

\tzpsection{Equation Grounding}
Equation anchors: proof=non sub chl sub or sub empirical registry sub definition sub combined,D: 0 sub Engine sub Outputspublish sub reassemble expanded sub context registry sub en; proof=candidate sub formal sub math registry sub definition sub combined,D: 0 sub Engine sub Outputspublish sub reassemble expanded sub context registry sub e; and proof=candidate sub formal sub math. Proof route: show that the composite map lands in the admissible target class once the TRAWIN ordering is respected. Formalization route: prove closure of the stated operator or functor under composition with the TRAWIN chain.

\tzpsection{Interpretive Role}
The entry converts a registry title into a transport map that can be reused by the surrounding transfer arc.
\endgroup

\tzpsection{TZP Thesaurus}
\begin{enumerate}[leftmargin=1.6em,itemsep=6pt,topsep=4pt]
\item \textbf{Maximum Coherence Targeting}: The algorithmic search method locating the specific spatial coordinate that maximizes the shared information between three bodies.
\item \textbf{Optimization of Entanglement Correlation}: The mathematical routine filtering out background noise to isolate the strongest three-way quantum linkage.
\item \textbf{Multi-Node State Distillation}: The protocol extracting a clean, readable signal from a highly distributed and overlapping informational network.
\end{enumerate}

\tzpsection{Encyclopedia Mathematica}
\begingroup\footnotesize\sloppy
The visual plate below is reserved for the source-specific equation and progression graphic for [ID 0349]. It is included only as a companion to the canonical equation, not as a substitute for the formal text.
\par\endgroup
\ifTZPDarkEdition
\tzpdictionaryvisualband{D:/TZPID/kdp_workflow/build/tikz_figures/ID0349_dictionary_figure_dark.pdf}
\fi

\clearpage
\begingroup\tzpsupportsetup
\noindent\textbf{\fontsize{10.8pt}{12.8pt}\selectfont [ID 0349] Constructive Logic and Proof Certificate}\par
\noindent\textbf{\fontsize{10pt}{12pt}\selectfont Tripartite Information Retrieval Algorithm}\par
\vspace{0.45em}
\begingroup
\footnotesize
\setlength{\parskip}{0.22em}
\setlength{\parindent}{0pt}
\noindent\textbf{Curry--Howard--Lambek Interpretation}\par
Under the Curry--Howard--Lambek correspondence, this registry entry is read as a typed proof
object: the stated source condition supplies the proposition, the derivation supplies the
morphism, and the proof certificate records the machine handles needed to check the obligation
in the formal registry.

\vspace{0.45em}
\noindent\textbf{CHL Proof}\par
\textbf{Statement:} to retrieve information across the tripartite channels: (1) measure sequences of
DAANSphere addresses corresponding to mediator channel events; (2)...

\textbf{Logic Validation:}\\
\textbf{Proposition:} ``Tripartite Information Retrieval Algorithm is valid as a spectral relation when
frequency, phase, and propagation variables satisfy the stated equation.''\\
\textbf{Proof:} The source statement describes a wave, harmonic, or resonance relation. The proof
obligation is that the frequency-domain quantities are typed consistently and that the
derived propagation or resonance condition follows from the stated variables.

\textbf{VERDICT:} \ensuremath{\checkmark}\ \textbf{TRUE} --- spectral-relation validation

\textbf{Formal Status:} \ensuremath{\checkmark}\ THEORETICAL -- source-backed CHL proof obligation

\textbf{Physical Status:} THEORETICAL -- requires model-specific empirical interpretation
\par\vspace{0.45em}
\noindent{\tiny\textbf{Isabelle/HOL.} \tzpplaincode{physics\_validity\_from\_model, external\_validity\_package, registry\_number\_matches, equation\_count\_matches}}\par
\ifTZPDarkEdition
\tzpproofvisualband{D:/TZPID/kdp_workflow/build/tikz_figures/ID0349_proof_certificate_figure_dark.pdf}
\fi
\endgroup
\endgroup
